% =====================================================================
%  Rapport de Projet - Deep Learning Avance - TEK-UP University
%  Sujet : DeepGuard - Detection de visages generes par IA (Deepfake)
%  Auteurs : Sourour Ben Salha & Ghaith Aouissaoui
%  Compilation : pdflatex main.tex (2 fois) -- ou Overleaf
% =====================================================================
\documentclass[11pt,a4paper]{article}

\usepackage[utf8]{inputenc}
\usepackage[T1]{fontenc}
\usepackage{lmodern}
\usepackage[french]{babel}
\usepackage[a4paper,top=2.5cm,bottom=2.5cm,left=2.5cm,right=2.5cm,headheight=14pt]{geometry}
\usepackage{microtype}
\usepackage{graphicx}
\usepackage{xcolor}
\usepackage{booktabs,tabularx,multirow}
\usepackage{caption}
\usepackage{enumitem}
\usepackage{fancyhdr}
\usepackage{titlesec}
\usepackage{listings}
\usepackage[most]{tcolorbox}
\usepackage{tikz}
\usetikzlibrary{arrows.meta,positioning,fit,backgrounds,calc,babel}
\usepackage[hidelinks]{hyperref}
\usepackage{float}

% ---------------- Couleurs (inspirees du logo TEK-UP) ----------------
\definecolor{primaire}{HTML}{1E2A78}
\definecolor{secondaire}{HTML}{2E6DB4}
\definecolor{accent}{HTML}{C0392B}
\definecolor{vert}{HTML}{1E8449}
\definecolor{orange}{HTML}{D35400}
\definecolor{violet}{HTML}{6C3483}
\definecolor{grisclair}{HTML}{F4F6F8}
\definecolor{grismoyen}{HTML}{BFC5CC}
\definecolor{codebg}{HTML}{F7F7F9}

% ---------------- Informations ----------------
\newcommand{\nomprojet}{DeepGuard}
\newcommand{\etudiantA}{Sourour \textsc{Ben Salha}}
\newcommand{\etudiantB}{Ghaith \textsc{Aouissaoui}}
\newcommand{\enseignant}{M. Haythem \textsc{Ghazouani}}

% ---------------- Mise en forme ----------------
\titleformat{\section}{\Large\sffamily\bfseries\color{primaire}}{\thesection}{0.8em}{}[{\color{primaire}\titlerule[0.6pt]}]
\titleformat{\subsection}{\large\sffamily\bfseries\color{secondaire}}{\thesubsection}{0.8em}{}
\titlespacing*{\section}{0pt}{18pt}{10pt}

\pagestyle{fancy}
\fancyhf{}
\fancyhead[L]{\small\sffamily\color{primaire}\textbf{Deep Learning Avancé} | \nomprojet}
\fancyhead[R]{\small\sffamily\color{gray}TEK-UP University}
\fancyfoot[L]{\small\sffamily\color{gray}\etudiantA{} \& \etudiantB}
\fancyfoot[R]{\small\sffamily\color{primaire}\thepage}
\renewcommand{\headrulewidth}{0.4pt}
\renewcommand{\footrulewidth}{0.4pt}

\setlist{itemsep=2pt,topsep=4pt}
\captionsetup{font=small,labelfont={bf,color=primaire}}
\setlength{\parskip}{5pt}
\setlength{\parindent}{0pt}

\lstset{backgroundcolor=\color{codebg},basicstyle=\ttfamily\footnotesize,
  frame=single,rulecolor=\color{grismoyen},breaklines=true,columns=fullflexible,
  keepspaces=true,showstringspaces=false}

\newtcolorbox{encadre}[2][]{enhanced,colback=grisclair,colframe=#2,
  fonttitle=\sffamily\bfseries,coltitle=white,arc=2mm,boxrule=0.8pt,#1}

\tikzset{
  bloc/.style={draw=#1,fill=#1!10,rounded corners=3pt,thick,align=center,
               font=\small\sffamily,minimum height=1cm,text width=2.6cm},
  bloc/.default=secondaire,
  fleche/.style={-{Stealth[length=2.5mm]},thick,color=#1},
  fleche/.default=primaire,
  etiquette/.style={font=\scriptsize\sffamily,color=#1,fill=white,inner sep=1pt},
  etiquette/.default=gray,
}

\begin{document}

% =====================================================================
%  PAGE DE GARDE
% =====================================================================
\begin{titlepage}
\thispagestyle{empty}
\begin{tikzpicture}[remember picture,overlay]
  \fill[primaire] (current page.north west) rectangle ([yshift=-0.6cm]current page.north east);
  \fill[secondaire] ([yshift=-0.6cm]current page.north west) rectangle ([yshift=-0.75cm]current page.north east);
  \fill[primaire] (current page.south west) rectangle ([yshift=0.6cm]current page.south east);
\end{tikzpicture}
\centering
\vspace*{0.3cm}

\includegraphics[width=7cm]{logo.png}

\vspace{0.8cm}
{\normalsize\sffamily Cycle Ingénieur en Intelligence Artificielle \& Data Science --- 3\textsuperscript{ème} année}

\vspace{1.8cm}
{\large\sffamily\color{secondaire} RAPPORT DE PROJET DE MODULE}\\[0.2cm]
{\Large\sffamily\bfseries Deep Learning Avancé}

\vspace{1.4cm}
% ---------------- Sujet entre deux barres ----------------
{\color{primaire}\rule{\linewidth}{1.6pt}}\\[0.5cm]
{\huge\sffamily\bfseries\color{primaire} \nomprojet}\\[0.35cm]
{\LARGE\sffamily\bfseries Détection de Visages Générés par IA\\[0.15cm] (Deepfake : Real vs Fake Faces)}\\[0.4cm]
{\color{primaire}\rule{\linewidth}{1.6pt}}

\vfill

\begin{minipage}[t]{0.48\textwidth}
  \raggedright
  {\sffamily\bfseries\color{primaire} Réalisé par :}\\[0.2cm]
  {\sffamily\large \etudiantA}\\[0.1cm]
  {\sffamily\large \etudiantB}
\end{minipage}%
\hfill
\begin{minipage}[t]{0.48\textwidth}
  \raggedleft
  {\sffamily\bfseries\color{primaire} Encadré par :}\\[0.2cm]
  {\sffamily\large \enseignant}
\end{minipage}

\vspace{1.6cm}
{\sffamily\large Année universitaire 2026 -- 2027}
\vspace{0.8cm}
\end{titlepage}

\tableofcontents
\newpage

% =====================================================================
\section{Introduction}
% =====================================================================

Les progrès de l'intelligence artificielle générative, en particulier des réseaux antagonistes
génératifs (GAN) comme StyleGAN, permettent aujourd'hui de créer des visages humains
photoréalistes de personnes qui n'existent pas. Ces images, appelées \emph{deepfakes}, sont
souvent impossibles à distinguer de vraies photographies à l'œil nu.

Ce projet, réalisé dans le cadre du module \emph{Deep Learning Avancé}, consiste à développer
\textbf{\nomprojet}, une application web capable de détecter automatiquement si un visage est
\textbf{réel} ou \textbf{généré par IA}, d'expliquer sa décision, et d'être déployée comme un
véritable produit logiciel en suivant les bonnes pratiques MLOps.

% =====================================================================
\section{Problématique}
% =====================================================================

La diffusion de visages synthétiques crée des risques concrets :
\begin{itemize}
  \item \textbf{Fraude à l'identité} : faux documents, contournement des vérifications en ligne ;
  \item \textbf{Faux profils} sur les réseaux sociaux, utilisés pour des arnaques ;
  \item \textbf{Désinformation} : création de faux témoins ou de fausses sources.
\end{itemize}

Un humain ne peut pas vérifier manuellement des milliers d'images, et les indices de
falsification (texture de la peau, arrière-plan, reflets dans les yeux, artefacts invisibles)
sont trop subtils.

\begin{encadre}[title={Problématique}]{primaire}
Comment concevoir un système de Deep Learning capable de \textbf{distinguer de façon fiable un
visage réel d'un visage généré par IA}, d'\textbf{expliquer visuellement} sa décision, et de le
\textbf{mettre à disposition} des utilisateurs à travers une application web rapide, fiable et
reproductible ?
\end{encadre}

% =====================================================================
\section{Objectifs}
% =====================================================================

\subsection{Objectif général}
Développer une application complète de détection de deepfakes, depuis la préparation des données
jusqu'au déploiement conteneurisé.

\subsection{Objectifs spécifiques}
\begin{enumerate}
  \item \textbf{Analyser les données} : explorer le jeu de données de visages réels et générés.
  \item \textbf{Entraîner et comparer des modèles} : une baseline CNN, un CNN pré-entraîné
        (EfficientNet) et un Vision Transformer (ViT).
  \item \textbf{Expliquer les prédictions} grâce à des cartes de chaleur (Grad-CAM, Attention Rollout).
  \item \textbf{Optimiser le modèle} en l'exportant au format ONNX pour une inférence rapide.
  \item \textbf{Développer une API REST} avec FastAPI pour servir le modèle.
  \item \textbf{Développer une interface web} avec React (upload d'image, webcam, affichage du résultat).
  \item \textbf{Industrialiser le projet} : versioning (DVC), suivi des expériences (MLflow),
        conteneurisation (Docker) et intégration continue (GitHub Actions).
\end{enumerate}

% =====================================================================
\section{Description du projet}
% =====================================================================

\subsection{Principe de fonctionnement}
L'utilisateur fournit une photo de visage à l'application (par upload, depuis une galerie
d'exemples ou via la webcam). L'image est envoyée à l'API, qui la prétraite et la transmet au
modèle de Deep Learning. Le modèle retourne :
\begin{itemize}
  \item le \textbf{verdict} : \og RÉEL \fg{} ou \og FAKE \fg{} ;
  \item la \textbf{probabilité} associée (niveau de confiance) ;
  \item une \textbf{carte de chaleur} qui montre les zones du visage ayant influencé la décision ;
  \item le \textbf{temps de réponse} en millisecondes.
\end{itemize}

\subsection{Jeux de données}

\begin{table}[H]
\centering
\caption{Jeux de données utilisés}
\small
\begin{tabularx}{\linewidth}{@{}lXX@{}}
\toprule
& \textbf{140k Real and Fake Faces} & \textbf{Real and Fake Face Detection} \\
\midrule
Source & Kaggle & Kaggle (Yonsei University) \\
Contenu & 70\,000 visages réels (FFHQ) + 70\,000 visages générés (StyleGAN) &
1\,081 visages réels + 960 visages retouchés \\
Résolution & $256\times256$ & $600\times600$ \\
Découpage & Train 100\,000 / Validation 20\,000 / Test 20\,000 & --- \\
Utilisation & Entraînement et évaluation & Test de généralisation \\
\bottomrule
\end{tabularx}
\end{table}

\subsection{Modèles de Deep Learning}

Le projet s'inscrit dans l'\textbf{Axe 4 : Vision Transformers}, avec une comparaison contre une
baseline CNN.

\begin{table}[H]
\centering
\caption{Modèles comparés}
\small
\begin{tabularx}{\linewidth}{@{}lX@{}}
\toprule
\textbf{Modèle} & \textbf{Description} \\
\midrule
Baseline CNN & Petit réseau convolutif de 4 blocs, entraîné from scratch. Sert de référence. \\
EfficientNet-V2-S & CNN pré-entraîné sur ImageNet, fine-tuné sur nos données (Transfer Learning). \\
ViT-B/16 & Vision Transformer : l'image est découpée en patchs de $16\times16$ pixels traités par
un mécanisme d'auto-attention, ce qui permet de comparer des zones éloignées du visage. \\
\bottomrule
\end{tabularx}
\end{table}

Pour rendre les décisions compréhensibles, nous utilisons \textbf{Grad-CAM} pour les CNN et
\textbf{Attention Rollout} pour le ViT.

\subsection{Fonctionnalités de l'application}
\begin{itemize}
  \item Upload d'image par glisser-déposer ;
  \item Galerie d'exemples testables en un clic ;
  \item Capture par webcam ;
  \item Affichage du verdict, d'une barre de probabilité et de la carte de chaleur ;
  \item Affichage de la latence côté serveur et côté client ;
  \item Documentation automatique de l'API (Swagger).
\end{itemize}

\subsection{Répartition des rôles}

\begin{table}[H]
\centering
\caption{Répartition des tâches dans le binôme}
\small
\begin{tabularx}{\linewidth}{@{}p{3.8cm}X@{}}
\toprule
\textbf{Membre} & \textbf{Responsabilités} \\
\midrule
\etudiantA & Deep Learning \& MLOps : données, entraînement, explicabilité, DVC, MLflow, export ONNX. \\
\etudiantB & Full-Stack \& DevOps : API FastAPI, interface React, Docker, CI/CD GitHub Actions. \\
\bottomrule
\end{tabularx}
\end{table}

% =====================================================================
\section{Technologies utilisées}
% =====================================================================

\begin{table}[H]
\centering
\caption{Stack technologique du projet}
\small
\begin{tabularx}{\linewidth}{@{}llX@{}}
\toprule
\textbf{Domaine} & \textbf{Technologie} & \textbf{Rôle dans le projet} \\
\midrule
\multirow{3}{*}{Deep Learning} & PyTorch & Construction et entraînement des modèles \\
 & timm & Modèles pré-entraînés (EfficientNet, ViT) \\
 & Grad-CAM & Explicabilité des prédictions \\
\midrule
\multirow{3}{*}{MLOps} & DVC & Versioning des données et des modèles \\
 & MLflow & Suivi des expériences et registre des modèles \\
 & ONNX Runtime & Optimisation et inférence rapide du modèle \\
\midrule
\multirow{2}{*}{Backend} & FastAPI & API REST asynchrone \\
 & Pydantic & Validation des données d'entrée et de sortie \\
\midrule
\multirow{2}{*}{Frontend} & React + Vite & Interface web interactive \\
 & TailwindCSS & Design moderne de l'interface \\
\midrule
\multirow{3}{*}{DevOps} & Docker & Conteneurisation des services \\
 & Docker Compose & Lancement de toute l'application en une commande \\
 & GitHub Actions & Intégration continue (tests et build automatiques) \\
\midrule
Tests & Pytest & Tests unitaires de l'API \\
\bottomrule
\end{tabularx}
\end{table}

% =====================================================================
\section{Architecture du projet}
% =====================================================================

\subsection{Architecture globale}

L'application est conçue comme un \textbf{système découplé} : chaque composant a un rôle unique
et communique avec les autres via des interfaces standard. L'ensemble est orchestré par Docker
Compose.

\begin{figure}[H]
\centering
\begin{tikzpicture}[every node/.style={font=\small\sffamily}]
  \node[bloc=primaire,text width=2.3cm] (user) {Utilisateur\\Navigateur web};
  \node[bloc=secondaire,text width=2.6cm,right=1.6cm of user] (front) {\textbf{Frontend}\\React\\(port 3000)};
  \node[bloc=secondaire,text width=2.6cm,right=1.6cm of front] (api) {\textbf{API REST}\\FastAPI + Pydantic\\(port 8000)};
  \node[bloc=vert,text width=2.6cm,below=1.5cm of api] (onnx) {\textbf{Modèle IA}\\ONNX Runtime\\+ Grad-CAM};
  \node[bloc=violet,text width=2.6cm,left=1.6cm of onnx] (mlflow) {\textbf{MLflow}\\Suivi \& registre\\(port 5000)};
  \node[bloc=orange,text width=2.3cm,left=1.6cm of mlflow] (dvc) {\textbf{DVC}\\Données \&\\modèles};
  \node[bloc=accent,text width=2.6cm,below=1.3cm of mlflow] (train) {\textbf{Entraînement}\\PyTorch (GPU)};

  \draw[fleche,{Stealth[length=2.5mm]}-{Stealth[length=2.5mm]}] (user) -- node[etiquette,above]{HTTP} (front);
  \draw[fleche,{Stealth[length=2.5mm]}-{Stealth[length=2.5mm]}] (front) -- node[etiquette,above]{JSON} (api);
  \draw[fleche] (api) -- node[etiquette,right]{image} (onnx);
  \draw[fleche=orange] (dvc) |- node[etiquette,pos=0.75,below]{données} (train);
  \draw[fleche=violet] (train) -- node[etiquette,right]{métriques} (mlflow);
  \draw[fleche=violet] (mlflow) -- node[etiquette,below]{meilleur modèle} (onnx);

  \begin{scope}[on background layer]
    \node[draw=primaire,dashed,rounded corners=6pt,fill=grisclair,inner sep=10pt,
          fit=(front)(api)(onnx)(mlflow),
          label={[font=\scriptsize\sffamily\bfseries,color=primaire]above right:Docker Compose}] {};
  \end{scope}
\end{tikzpicture}
\caption{Architecture globale de \nomprojet}
\end{figure}

\begin{table}[H]
\centering
\caption{Rôle de chaque composant}
\small
\begin{tabularx}{\linewidth}{@{}lX@{}}
\toprule
\textbf{Composant} & \textbf{Rôle} \\
\midrule
Frontend React & Interface utilisateur : envoi de l'image et affichage du résultat. \\
API FastAPI & Reçoit l'image, la valide, appelle le modèle et renvoie la réponse en JSON. \\
Modèle ONNX & Calcule la prédiction et la carte de chaleur. \\
MLflow & Enregistre les expériences et conserve le meilleur modèle. \\
DVC & Versionne les données et les poids des modèles en dehors de Git. \\
\bottomrule
\end{tabularx}
\end{table}

\subsection{Pipeline du projet (MLOps)}

\begin{figure}[H]
\centering
\begin{tikzpicture}[node distance=0.55cm, every node/.style={font=\scriptsize\sffamily}]
  \node[bloc=orange,text width=2.1cm] (d1) {\textbf{Données}\\Kaggle};
  \node[bloc=orange,text width=2.1cm,right=of d1] (d2) {\textbf{DVC}\\versioning};
  \node[bloc=secondaire,text width=2.1cm,right=of d2] (d3) {\textbf{Préparation}\\nettoyage, split};
  \node[bloc=accent,text width=2.1cm,right=of d3] (d4) {\textbf{Entraînement}\\PyTorch};
  \node[bloc=violet,text width=2.1cm,right=of d4] (d5) {\textbf{MLflow}\\suivi des runs};
  \node[bloc=vert,text width=2.1cm,below=1.1cm of d5] (d6) {\textbf{Export}\\ONNX};
  \node[bloc=primaire,text width=2.1cm,left=of d6] (d7) {\textbf{API}\\FastAPI};
  \node[bloc=primaire,text width=2.1cm,left=of d7] (d8) {\textbf{Interface}\\React};
  \node[bloc=gray,text width=2.1cm,left=of d8] (d9) {\textbf{Docker}\\Compose};
  \node[bloc=gray,text width=2.1cm,left=of d9] (d10) {\textbf{CI/CD}\\GitHub Actions};
  \draw[fleche] (d1) -- (d2); \draw[fleche] (d2) -- (d3); \draw[fleche] (d3) -- (d4);
  \draw[fleche] (d4) -- (d5); \draw[fleche] (d5) -- (d6); \draw[fleche] (d6) -- (d7);
  \draw[fleche] (d7) -- (d8); \draw[fleche] (d8) -- (d9); \draw[fleche] (d9) -- (d10);
\end{tikzpicture}
\caption{Pipeline de bout en bout : des données au déploiement}
\end{figure}

\subsection{Déroulement d'une prédiction}

\begin{figure}[H]
\centering
\begin{tikzpicture}[every node/.style={font=\scriptsize\sffamily}]
  \foreach \x/\nom/\c in {0/Utilisateur/primaire, 3.6/React/secondaire, 7.2/FastAPI/secondaire, 10.8/Modèle ONNX/vert} {
    \node[bloc=\c,text width=2.2cm,minimum height=0.7cm] at (\x,0) {\textbf{\nom}};
    \draw[gray,dashed] (\x,-0.4) -- (\x,-5.4);
  }
  \draw[fleche] (0,-1.0) -- node[etiquette,above]{1. choisit une image} (3.6,-1.0);
  \draw[fleche] (3.6,-1.8) -- node[etiquette,above]{2. POST /predict} (7.2,-1.8);
  \draw[fleche] (7.2,-2.6) -- node[etiquette,above]{3. image prétraitée} (10.8,-2.6);
  \draw[fleche=vert] (10.8,-3.4) -- node[etiquette,above]{4. probabilités + carte} (7.2,-3.4);
  \draw[fleche] (7.2,-4.2) -- node[etiquette,above]{5. réponse JSON} (3.6,-4.2);
  \draw[fleche] (3.6,-5.0) -- node[etiquette,above]{6. verdict affiché} (0,-5.0);
\end{tikzpicture}
\caption{Échanges entre les composants lors d'une prédiction}
\end{figure}

\subsection{Endpoints de l'API}

\begin{table}[H]
\centering
\caption{Routes exposées par l'API FastAPI}
\small
\begin{tabularx}{\linewidth}{@{}llX@{}}
\toprule
\textbf{Méthode} & \textbf{Route} & \textbf{Description} \\
\midrule
GET & \texttt{/health} & Vérifie que le serveur et le modèle sont prêts \\
GET & \texttt{/model/metadata} & Informations sur le modèle (architecture, classes, version) \\
POST & \texttt{/predict} & Reçoit une image et renvoie le verdict et la probabilité \\
POST & \texttt{/predict/visualize} & Renvoie le verdict et un calque de chaleur (PNG) à superposer à la photo \\
\bottomrule
\end{tabularx}
\end{table}

\subsection{Organisation du dépôt Git}

\begin{lstlisting}
deepguard/
|-- .github/workflows/   # CI/CD (tests + build Docker)
|-- data/                # Donnees suivies par DVC
|-- src/                 # Code Deep Learning (dataset, modeles, entrainement)
|-- backend/             # API FastAPI + tests Pytest + Dockerfile
|-- frontend/            # Application React + Dockerfile
|-- docker-compose.yml   # Lancement de toute l'application
|-- dvc.yaml             # Pipeline reproductible
`-- README.md            # Guide d'installation
\end{lstlisting}

L'application complète se lance avec une seule commande :
\texttt{docker compose up --build}.

% =====================================================================
\section{Conclusion}
% =====================================================================

Ce projet propose \nomprojet, une solution complète de détection de visages générés par IA. Il
combine des modèles de Deep Learning avancés (CNN et Vision Transformer), des techniques
d'explicabilité, et une architecture logicielle professionnelle (FastAPI, React, Docker) avec un
cycle de vie MLOps (DVC, MLflow, CI/CD).

Comme perspectives, nous envisageons d'entraîner le modèle sur des images issues de modèles de
diffusion récents, d'étendre la détection aux vidéos, et de déployer l'application sur le cloud.

\end{document}
